% Options for packages loaded elsewhere
\PassOptionsToPackage{unicode}{hyperref}
\PassOptionsToPackage{hyphens}{url}
\documentclass[
  11pt,
]{article}
\usepackage{xcolor}
\usepackage[margin=1in]{geometry}
\usepackage{amsmath,amssymb}
\setcounter{secnumdepth}{-\maxdimen} % remove section numbering
\usepackage{iftex}
\ifPDFTeX
  \usepackage[T1]{fontenc}
  \usepackage[utf8]{inputenc}
  \usepackage{textcomp} % provide euro and other symbols
\else % if luatex or xetex
  \usepackage{unicode-math} % this also loads fontspec
  \defaultfontfeatures{Scale=MatchLowercase}
  \defaultfontfeatures[\rmfamily]{Ligatures=TeX,Scale=1}
\fi
\usepackage{lmodern}
\ifPDFTeX\else
  % xetex/luatex font selection
\fi
% Use upquote if available, for straight quotes in verbatim environments
\IfFileExists{upquote.sty}{\usepackage{upquote}}{}
\IfFileExists{microtype.sty}{% use microtype if available
  \usepackage[]{microtype}
  \UseMicrotypeSet[protrusion]{basicmath} % disable protrusion for tt fonts
}{}
\makeatletter
\@ifundefined{KOMAClassName}{% if non-KOMA class
  \IfFileExists{parskip.sty}{%
    \usepackage{parskip}
  }{% else
    \setlength{\parindent}{0pt}
    \setlength{\parskip}{6pt plus 2pt minus 1pt}}
}{% if KOMA class
  \KOMAoptions{parskip=half}}
\makeatother
\setlength{\emergencystretch}{3em} % prevent overfull lines
\providecommand{\tightlist}{%
  \setlength{\itemsep}{0pt}\setlength{\parskip}{0pt}}
\usepackage{bookmark}
\IfFileExists{xurl.sty}{\usepackage{xurl}}{} % add URL line breaks if available
\urlstyle{same}
% fallback for those not using the hyperref driver hyperxmp:
\makeatletter
\@ifundefined{xmpquote}{\newcommand{\xmpquote}[1]{#1}}{}
\makeatother
\hypersetup{
  pdftitle={Simultaneous avoidance of countable families},
  hidelinks,
  pdfcreator={LaTeX via pandoc}}

\title{Simultaneous avoidance of countable families}
\author{}
\date{21 September 2026 --- continuation for discussion}

\begin{document}
\maketitle

\textbf{Local theorem, with a hand proof below.} A countable family of
sequences over a fixed finite alphabet has a single avoiding population
if and only if none of its members is eventually periodic. For an
alphabet of size \(n\ge2\), the population can have \(n\) roots and at
most \(n\) children per vertex, which is the smallest possible uniform
outdegree bound. If the sequences have a uniformly computable
enumeration and are all non-eventually-periodic, such a population can
be constructed computably. Independent review and external priority
checking remain; this addition has not been checked in Lean.

The earlier explicit \(P_s\) construction avoided one prescribed
sequence; \texttt{complement-and-barriers.md} proves that it also avoids
its complement. The present result treats arbitrary countable families,
using a different label sequence built by successive finite extensions.

\section{1. A common space of
populations}\label{a-common-space-of-populations}

For any \(r\in\{0,1\}^{\mathbb N}\), let \(G_r\) have vertices
\(\mathbb N\), birthdate \(v\) at \(v\), roots 0 and 1, and precisely
these incoming edges for each \(v\ge2\):

\[v-1\longrightarrow v\text{ labelled }r(v),\qquad
v-2\longrightarrow v\text{ labelled }1-r(v).\]

There is no edge \(0\to1\). Every \(G_r\) satisfies the population
axioms: infinite vertex set, finitely many roots, finite birthdate
sublevels, birthdates increasing along edges, one incoming parent of
each label at every non-root, and outdegree at most two. Every infinite
path has increments one or two. The values \(r(0),r(1)\) do not label
edges but are harmless coordinates.

\section{2. Finite-extension lemma}\label{finite-extension-lemma}

\textbf{Lemma.} Fix a non-eventually-periodic binary sequence \(s\) and
any prescribed prefix \(r(0),\ldots,r(N)\), with \(N\ge1\). There is an
extension \(r\) for which \(G_r\) avoids \(s\) from every starting
vertex.

Set \(C=-N-2\) and prescribe the unassigned tail by

\[r(v)=s\!\left(\left\lfloor\frac{v-C}{2}\right\rfloor\right)
\mathbin{\mathrm{xor}}((v-C)\bmod2),\qquad v>N.\]

All indices on the right are nonnegative. The old prefix is unchanged.
Suppose there were an infinite \(s\)-matching path \(v_0,v_1,\ldots\)
and put \(d_k=v_k-2k\). Each step changes \(d_k\) by \(-1\) or 0.

Let \(j\) be its first index with \(v_j>N\); this exists because
vertices strictly increase. If \(j=0\), then \(d_j=v_0\ge0>C\).
Otherwise \(v_j\le N+2\) because edges have length at most two. Also
\(j\le v_j\) because \(v_0\ge0\) and every increment is at least one.
Consequently

\[d_j=v_j-2j\ge-v_j\ge-N-2=C.\]

For \(k\ge j\), the lower bound \(d_k\ge C\) is invariant. At equality,
\(v_k=2k+C\) and an increment-one edge would have label

\[r(v_k+1)=s(k)\mathbin{\mathrm{xor}}1,\]

so it cannot match. If \(d_k>C\), decreasing by at most one preserves
the bound. The edge endpoints used here are beyond \(N\), so the
prescribed tail formula really applies; the arbitrary old prefix is no
longer involved.

The bounded nonincreasing integer \(d_k\) therefore stabilizes at some
\(D\ge C\). All subsequent increments are two. Write \(h=D-C\ge0\) and
\(q=\lfloor h/2\rfloor+1>0\). Matching the increment-two edge says

\[s(k)=1-r(2k+D+2).\]

For even \(h\) this is \(s(k)=1-s(k+q)\); for odd \(h\) it is
\(s(k)=s(k+q)\). These give eventual periods \(2q\) and \(q\),
respectively, contradicting the hypothesis on \(s\). This proves the
lemma, including arbitrary starts and arbitrary preserved finite
prefixes.

\section{3. Simultaneous avoidance}\label{simultaneous-avoidance}

For a fixed target \(s\) and start \(i\), define

\[U_{s,i}=\{r:G_r\text{ has no infinite }s\text{-matching path from }i\}.\]

This set is open in Cantor space. The matching-prefix tree is finitely
branching. By König's lemma, absence of an infinite branch means that
some finite matching level is empty. A level of depth \(b\) only reads
labels at vertices at most \(i+2b\), so its emptiness is preserved by
fixing a finite prefix of \(r\).

If \(s\) is non-eventually-periodic, \(U_{s,i}\) is dense: any finite
cylinder can first be extended through coordinate 1, then extended by
the preceding lemma to a member of \(U_{s,i}\).

Now let \((s_j)_{j\in J}\) be a countable family of
non-eventually-periodic binary sequences. The set of simultaneous
avoiders is

\[\bigcap_{j\in J}\bigcap_{i\in\mathbb N}U_{s_j,i}.\]

It is a countable intersection of dense open sets in Cantor space, hence
comeagre and nonempty by the Baire category theorem. Any member gives
the required single two-root, degree-two population. An empty family is
trivial.

Conversely, a family containing an eventually periodic member cannot
have an avoiding population, by Alexander's eventual-periodic
unavoidability theorem (with the boundary repair recorded in
\texttt{results.md}). This proves the stated classification for
countable binary families.

This is a category statement about the choice of \textbf{edge labels
\(r\)}, not an almost-sure statement under random labels. No measure-one
claim is made.

\section{4. Effective construction under a uniform
promise}\label{effective-construction-under-a-uniform-promise}

The Baire proof can be implemented directly. Assume \(s_j(k)\) is
uniformly computable in \((j,k)\), with the promise that every \(s_j\)
is non-eventually periodic. Enumerate all requirements \((j,i)\), and
maintain a finite prefix of \(r\) extending through at least coordinate
1.

At a stage for \((j,i)\):

\begin{enumerate}
\def\labelenumi{\arabic{enumi}.}
\tightlist
\item
  Form the lemma's temporary infinite extension of the current prefix
  using target \(s_j\).
\item
  Enumerate the finite matching levels from \(i\) in this temporary
  population until the first empty level. The lemma and König's lemma
  guarantee that this search terminates.
\item
  Freeze all bits of that extension read by the search, together with at
  least one further coordinate beyond the old prefix.
\end{enumerate}

The finite barrier for this requirement persists under every later
extension. Every requirement is visited, and the prefix grows at every
stage. To compute one desired output bit, run stages until that bit has
been frozen. Each stage terminates and uses finitely many values of the
uniformly computable targets, so the final \(r\) is computable. The same
construction works relative to an oracle presenting an arbitrary
sequence of targets.

No algorithm recognizing the promise is asserted. There is no
contradiction with taking the countable family of \textbf{all computable
aperiodic sequences}: the existence theorem handles that family, but it
has no uniformly computable enumeration containing only aperiodic
sequences and covering them all. Indeed, such an enumeration would yield
a computable \(r\) avoiding all of them. If \(r\) were aperiodic, the
increment-one path \(1,2,3,\ldots\) would realize the computable
aperiodic sequence \(r(2),r(3),\ldots\), a contradiction. If \(r\) were
eventually periodic, that same path realizes an eventually periodic word
but does not by itself settle the contradiction; use instead the next
paragraph's periodic-tail argument.

If \(r\) is eventually periodic with period \(p\), \(G_r\) realizes
every binary sequence. For any desired finite word, choose its endpoint
sufficiently far past the exceptional prefix and follow its uniquely
labelled incoming parents backwards. All vertices of this finite path
can be kept in the periodic tail. Translate the path backwards by a
multiple of \(p\) until its start lies in a fixed tail interval of \(p\)
vertices. The edge labels remain the same and all vertices remain in the
tail. Arbitrarily long prescribed prefixes are thus realized from this
finite set of starts. König's lemma gives an infinite realizing path.
This proves universality and completes the contradiction for eventually
periodic \(r\) as well.

\section{5. Optimal degree for every finite
alphabet}\label{optimal-degree-for-every-finite-alphabet}

Identify an alphabet of size \(n\ge2\) with \(\mathbb Z/n\mathbb Z\).
Use vertices \(\mathbb N\), roots \(0,\ldots,n-1\), birthdate \(v\),
and, for each \(v\ge n\), incoming edges \(v-a\to v\), \(1\le a\le n\).
At each endpoint assign the \(n\) labels bijectively to these edges. All
such assignments give populations of outdegree at most \(n\). The
parameter space is a product of finite permutation spaces, one for each
endpoint, and is a Baire space.

\textbf{General finite-extension lemma.} Preserve arbitrary
incoming-label permutations at endpoints through \(N\ge n-1\). For a
fixed aperiodic target \(s\), set \(C=-(n-1)(N+n)\) and label all later
edges by

\[\ell(v-a,v)=s\!\left(\left\lfloor\frac{v-C}{n}\right\rfloor\right)
+a-((v-C)\bmod n)-1\pmod n,\qquad v>N.\]

As \(a\) ranges from 1 to \(n\) these labels are a permutation, as
required. Suppose \(v_k\) were an infinite matching path and put
\(d_k=v_k-nk\). Its increments lie in \(\{1-n,\ldots,0\}\). At the first
entry \(j\) into vertices beyond \(N\), either \(j=0\) and \(d_j\ge0\),
or \(v_j\le N+n\) and \(j\le v_j\), giving

\[d_j\ge-(n-1)v_j\ge-(n-1)(N+n)=C.\]

To check that \(d_k\ge C\) persists, set \(h=d_k-C\ge0\). A step of
length \(a\) could cross below \(C\) only if \(h+a<n\). In that case
\(0\le h\le n-2\), the endpoint is \(v_{k+1}=nk+C+h+a\), and its edge
label is

\[s(k)+a-(h+a)-1=s(k)-h-1\pmod n.\]

Since \(1\le h+1\le n-1\), this differs from \(s(k)\). Such a crossing
cannot be matching. All relevant endpoints are beyond \(N\), so the tail
formula applies.

Thus \(d_k\) eventually stabilizes at \(D\ge C\) and every later edge
has length \(n\). Write \(h=D-C\), \(b=h\bmod n\) and
\(q=\lfloor h/n\rfloor+1>0\). Matching then requires

\[s(k)=s(k+q)+n-b-1\pmod n\]

for all sufficiently large \(k\). Iterating \(n\) times gives the
positive eventual period \(nq\), a contradiction. This proves the lemma.

Absence of a matching path from a fixed start is again open (depth \(b\)
reads only endpoints through \(i+nb\)) and dense by this lemma. The same
countable intersection and staged freezing arguments prove simultaneous
avoidance and its effective version with \(n\) roots and maximum
outdegree \(n\).

For a single target there is an even simpler formula: use the displayed
label rule with \(C=0\) at \emph{every} endpoint \(v\ge n\). The
invariant starts at \(v_0\ge0\), and the same proof applies immediately.
This generalizes the original binary construction directly, without a
binary code or vertex-copy degree overhead.

\textbf{Optimality of the degree bound.} In any \(n\)-label population
with \(R\) roots and maximum outdegree at most \(n-1\), take a finite
birthdate sublevel set \(F\) containing all roots and \(M\) vertices. It
is ancestor-closed. The incoming-label requirement gives at least
\(n(M-R)\) internal edges, while the outdegree assumption gives at most
\((n-1)M\). Hence \(M\le nR\). Birthdate sublevels have unbounded size
in an infinite population, a contradiction. Thus no such population
exists below degree \(n\) at all. This counting obstruction was already
recorded in the earlier supplement; the new point is attainment for
arbitrary countable aperiodic target families. The one-symbol case has
no aperiodic targets; an empty family admits a chain.

\section{6. What is now settled, and what is
not}\label{what-is-now-settled-and-what-is-not}

\begin{itemize}
\tightlist
\item
  Arbitrary finite collections of aperiodic targets over a fixed finite
  alphabet can be avoided simultaneously; the result extends to
  countable collections. The smallest possible uniform outdegree is
  attained for every finite alphabet.
\item
  In particular one population can avoid every shift and complement of
  Thue--Morse. It must nevertheless realize some other member of the
  Thue--Morse orbit closure, by the earlier subshift-intersection
  theorem.
\item
  There is a population avoiding every computable aperiodic binary
  sequence. For the present \(G_r\) family, its labelling cannot be
  computable, as just shown. No claim about every other computable
  population presentation is made.
\item
  All uncountable families cannot be handled in this way: the family of
  all aperiodic binary sequences is unavoidable as a set, since every
  binary population realizes continuum many aperiodic words.
  Countability is a sufficient hypothesis, not a proved necessary one
  for joint avoidance.
\end{itemize}

The remaining questions include a characterization of jointly avoidable
uncountable families, quantitative costs of the effective
finite-extension construction, sharp Thue--Morse matching barriers, and
the other embedding/ordinal variants in the original packet. None is
silently resolved by the countable intersection.

Finite controls in \texttt{evidence/44u-countable-avoidance-checks.py}
check arbitrary short preserved prefixes, the shifted-offset boundary,
and a small finite initial segment of the staged construction. They
cannot certify countably many requirements; that conclusion uses the
proofs above.

Source: Samuel A. Alexander, \emph{Biologically unavoidable sequences},
Electronic Journal of Combinatorics \textbf{20}(1) (2013), P31,
\href{https://doi.org/10.37236/3035}{doi:10.37236/3035}. The
finite-extension lemma and its use here are local continuations of the
earlier proposed classification proof. Baire category and König's lemma
are standard antecedents, not new results claimed by this note.

\end{document}
